\section{教育背景}

\textbf{电子科技大学 (UESTC) 通信工程学士} \hfill 2022年9月 --- 2026年6月

\textbf{格拉斯哥大学 (UofG) 电子信息工程学士 (双学位项目)} \hfill 2022年9月 --- 2026年6月

\textbf{香港中文大学 (深圳) 计算机科学 (在读)} \hfill 2026年9月 --- 2028年6月


\section{研究与项目经历}

\MYhref{https://github.com/Marcobisky/llmlab}{\textbf{LLMlab: 基于形式语言的 LLM RLVR 快速后训练验证平台}} \hfill 2026年6月

\begin{itemize}
    \item 提出动机: 为了降低 LLM 的极高算法验证成本, LLMlab 可在可控的合成数据上进行高效的算法性能验证.
    \item 设计了一种确定性的形式语言, 支持在不同形式语言难度上的 RLVR 的高效算法性能验证.
    \item 实现了在 2.67M、0.15M 的教师和学生模型上的完整 pretrain, SFT, GRPO, KD, OPD, SDPO 流水线和平行消融实验, 最高实现了在语言难度 0-3 等级的 100\% 正确率和难度 4-6 等级超过 70\% 的正确率.
    \item 集成了 PCA 投影的 loss landscape 及权重变化轨迹、逐层注意力热力图、exposure bias 度量等可视化工具集以便对模型的可解释性和训练过程进行深入分析.
\end{itemize}

\textbf{边缘 TinyML 的 RISCV 加速器系统级软硬件协同设计} \hfill 2025年9月 --- 2026年4月

\vspace{-0.2em} \begin{small} \textit{研究助理, \MYhref{https://scholar.google.com/citations?user=1NT1jFMAAAAJ&hl=en}{李耘 教授}, 电子科技大学} \end{small} \vspace{-0.2em}

\begin{itemize}
    \item 在 Artix-7 FPGA 上独立架构并开发了用于实时 YOLOv8n 边缘推理的神经网络处理单元 (NPU), 通过自定义 RISC-V 指令集扩展 (Xnpu) 实现了无软核处理器的纯硬件执行架构.
    \item 自研端到端 Python 机器学习编译器，实现了全自动 INT16 量化与内存感知的指令调度，将量化后的 mAP 精度损失严格控制在 PyTorch FP32 基线的 0.3\% 以内.
    \item 设计了基于 3×3 脉动阵列 (Systolic Array) 的参数化 RTL 算子，并实现了全硬件加速的后处理引擎 (DFL 框解码与 NMS)，峰值算力达 288 MACs/cycle 与 23.4 GOPS.
    \item 构建了完整的边缘计算端到端数据流，集成了异步摄像头/UDP 视频流管道及 AXI4 DDR3L 内存总线复用，并基于 \texttt{Cocotb}、\texttt{Bazel} 和 \texttt{Icarus Verilog} 完成了全系统的仿真与验证.
\end{itemize}

\MYhref{https://github.com/Marcobisky/CNN-for-mbed}{\textbf{面向嵌入式系统的 CNN/LSTM}} \hfill 2024年2月 --- 2024年5月

\begin{itemize}
    \item 设计并基于 STM32 微控制器实现了 CNN 与 LSTM 模型的端侧推理, 实现了高精度和低延迟的患者跌倒检测预警、体温监测及数据曲线的可视化.
    \item 本地化采集、标注并构建了三轴加速度时序数据集. 在 Linux 环境下完成模型训练, 使用 \texttt{C++} 和 MbedOS 库在 STM32 上进行了硬编码实现与算法加速.
\end{itemize}

\textbf{人声识别遥控智能小车} \hfill 2023年9月 --- 2023年12月

\begin{itemize}
    \item 领导一个4人团队, 使用 \texttt{C} 标准库在 STM32F103 微处理器上设计并实现了一辆声控小车, 支持前进/后退, 左转/右转, 左滑/右滑, 睡眠/关机等指令.
    \item 使用 MIT App Inventor 开发了一个图形化 Android 手机应用, 支持手机端蓝牙发送运动指令控制小车运动.
\end{itemize}


\section{相关技能}

\begin{tabularx}{\linewidth}{@{}l X@{}}
\textbf{编程语言} &  \normalsize{\texttt{Python}, \texttt{PyTorch}, \texttt{C/C++} , \texttt{Makefile}.}\\
\textbf{英语} &  \normalsize{英语四六级 (615, 573), \MYhref{https://marcobisky.github.io/cv/ielts-score.pdf}{雅思 7} (8/7/7/6.5)}\\
\textbf{音乐} &  \normalsize{二胡, 小提琴, 钢琴, Sibelius, MuseScore, ACE Studio, IMSLP, \MYhref{https://space.bilibili.com/479299347?spm_id_from=333.1007.0.0}{音乐自媒体}.}\\
\end{tabularx}


\section{奖项}

\textbf{电子科技大学一等学业奖学金 (前 5\%)} \hfill 2023年12月, 2024年12月 \\
\textbf{国家奖学金 (前 0.2\%)} \hfill 2024年12月 \\
\textbf{\MYhref{https://www.gla.uestc.edu.cn/info/1089/17553.htm}{2026 届省级优秀毕业生}} \hfill 2025年10月


\section{音乐经历}

\textbf{\MYhref{https://www.musiceol.com/news/html/2013-10/201310812184665677010.html}{中国音乐小金钟奖第二届全国二胡比赛优秀奖}} \hfill 2013年10月 \\
\textbf{南昌市青少年民族乐团二胡首席} \hfill 2016年9月 --- 2018年6月 \\
\textbf{\MYhref{http://www.moe.gov.cn/srcsite/A17/moe_794/moe_628/202408/t20240806_1144389.html}{全国第七届大学生艺术展演一等奖}} \hfill 2024年9月 \\
\textbf{电子科技大学交响乐团第二小提琴首席} \hfill 2024年2月 --- 2026年6月

